%==========================================================
%  CHAPITRE 1 — PRÉSENTATION DE L'ORGANISME D'ACCUEIL
%==========================================================
\chapter{Présentation de l'organisme d'accueil}
\label{chap:organisme}

\begin{remarque}
  Ce chapitre dresse un portrait de \textbf{\CompanyName},
  en présentant ses activités, son positionnement, son
  organisation interne et le département au sein duquel
  le stage s'est déroulé.
\end{remarque}

%──────────────────────────────────────────────────────────
\section{Présentation générale de l'entreprise}
\label{sec:org:presentation}
%──────────────────────────────────────────────────────────

\subsection{Fiche d'identité}

\begin{center}
\begin{tabular}{>{\bfseries\sffamily\color{SecGrey}}l l}
  \toprule
  Raison sociale       & \CompanyName \\
  Secteur d'activité   & \todo{ex : Informatique / Finance / Industrie} \\
  Forme juridique      & \todo{ex : SA, SARL, SAS, Établissement public} \\
  Date de création     & \todo{ex : 1995} \\
  Siège social         & \todo{Adresse complète} \\
  Site Web             & \url{\todo{https://www.entreprise.com}} \\
  Effectif             & \todo{ex : 250 employés} \\
  Chiffre d'affaires   & \todo{ex : XX M€ (20XX)} \\
  \bottomrule
\end{tabular}
\end{center}

\subsection{Activités et domaines d'expertise}

\todo{Décrire en 2--3 paragraphes les activités principales de
l'entreprise, ses clients, ses marchés et sa vision stratégique.
Exemple : « \CompanyName\ est spécialisée dans \ldots\ Elle propose
une gamme de solutions couvrant \ldots\ et sert une clientèle de
\ldots »}

\subsection{Présence géographique et chiffres clés}

\todo{Mentionner les implantations (nationales ou internationales),
les certifications (ISO, etc.), les partenariats ou labels obtenus.}

%──────────────────────────────────────────────────────────
\section{Organigramme et organisation interne}
\label{sec:org:organigramme}
%──────────────────────────────────────────────────────────

La figure~\ref{fig:organigramme} présente l'organigramme général
de \textbf{\CompanyName}.

\begin{figure}[H]
  \centering
  %─── Organigramme TikZ ────────────────────────────────
  \begin{tikzpicture}[
    node distance=1.5cm and 2cm,
    box/.style={
      rectangle, rounded corners=4pt,
      fill=ChapBlue!12, draw=ChapBlue, text=ChapBlue,
      font=\sffamily\small\bfseries,
      minimum width=3.2cm, minimum height=0.9cm,
      align=center, drop shadow
    },
    line/.style={draw=ChapBlue!60, -{Stealth[length=6pt]}, thick}
  ]
    % Direction
    \node[box, fill=ChapBlue, text=white] (DG) {Direction Générale};

    % Niveau 2
    \node[box, below left=of DG, xshift=-1cm]  (DIR1) {\todo{Direction 1}};
    \node[box, below=of DG]                    (DIR2) {\todo{Direction 2}};
    \node[box, below right=of DG, xshift=1cm]  (DIR3) {\todo{Direction 3}};

    % Niveau 3 (exemple sous DIR2)
    \node[box, below left=of DIR2]   (D21) {\todo{Dépt. A}};
    \node[box, below right=of DIR2]  (D22) {\todo{Dépt. B}};

    % Flèches
    \draw[line] (DG) -- (DIR1);
    \draw[line] (DG) -- (DIR2);
    \draw[line] (DG) -- (DIR3);
    \draw[line] (DIR2) -- (D21);
    \draw[line] (DIR2) -- (D22);
  \end{tikzpicture}
  \caption{Organigramme simplifié de \CompanyName\ \todo{(à adapter)}}
  \label{fig:organigramme}
\end{figure}

\todo{Commenter l'organigramme en 3--5 phrases : nombre de directions,
hiérarchie, place du département d'accueil.}

%──────────────────────────────────────────────────────────
\section{Département d'accueil et équipe}
\label{sec:org:departement}
%──────────────────────────────────────────────────────────

Le stage s'est déroulé au sein du département
\textbf{\todo{Nom du département}}, rattaché à la direction
\textbf{\todo{Nom de la direction}}.

\subsection{Missions du département}

\begin{itemize}
  \item \todo{Mission principale du département}
  \item \todo{Mission secondaire 1}
  \item \todo{Mission secondaire 2}
  \item \todo{Outils et systèmes gérés par le département}
\end{itemize}

\subsection{Composition de l'équipe}

\begin{center}
\begin{tabular}{>{\sffamily\bfseries}l l l}
  \toprule
  \textbf{Nom / Prénom} & \textbf{Rôle} & \textbf{Lien avec le stage} \\
  \midrule
  \SupervisorPro         & \todo{Titre/Poste}    & Encadrant direct \\
  \todo{Collègue 1}      & \todo{Développeur}    & Collaboration technique \\
  \todo{Collègue 2}      & \todo{Analyste}       & Revue fonctionnelle \\
  \todo{Collègue 3}      & \todo{DevOps}         & Déploiement \\
  \bottomrule
\end{tabular}
\end{center}

%──────────────────────────────────────────────────────────
\section{Environnement de travail}
\label{sec:org:environnement}
%──────────────────────────────────────────────────────────

\todo{Décrire brièvement les conditions de travail : locaux, horaires,
outils mis à disposition (matériel, licences, VPN, accès serveurs),
culture d'entreprise.}

\begin{itemize}
  \item \textbf{Infrastructure :} \todo{ex : Serveurs on-premise + Cloud AWS}
  \item \textbf{Outils collaboratifs :} \todo{ex : Slack, Teams, Confluence}
  \item \textbf{Environnement de développement :} \todo{ex : Linux Ubuntu 22.04 LTS}
  \item \textbf{Méthodologie de travail :} \todo{ex : Agile Scrum, sprints de 2 semaines}
\end{itemize}

\clearpage
